% ============================================================
% f2012_hubei_science_q04_views.tex
% 2012 湖北卷（理科）第 4 题 几何体三视图
% 正视图：直角梯形，底宽 2，左高 2，右高 4；侧视图：2×4 矩形内切半径 1 的圆；
% 俯视图：半径 1 的圆。几何体为半径 1、高 4 的圆柱被过轴心高度的斜平面所截。
% 来源 PDF：普通高考/2012/2012湖北理.pdf（SHA-256 b84a88bc…07cc）
% 原页：第 1 页（A3 横排）第 1 栏，印面 737，第 4 题题干下方
% ============================================================
\begin{tikzpicture}[>=Stealth, line width=0.7pt, font=\small, x=1cm, y=1cm]
  % ---------- 正视图 ----------
  \draw (0,0) -- (2,0) -- (2,4) -- (0,2) -- cycle;
  % 底宽 2
  \draw[<->] (0,-0.45) -- (2,-0.45);
  \draw (0,-0.58) -- (0,-0.32);
  \draw (2,-0.58) -- (2,-0.32);
  \node at (1,-0.85) {$2$};
  % 左高 2
  \draw[<->] (-0.45,0) -- (-0.45,2);
  \draw (-0.58,0) -- (-0.32,0);
  \draw (-0.58,2) -- (-0.32,2);
  \node at (-0.85,1) {$2$};
  % 右高 4
  \draw[<->] (2.45,0) -- (2.45,4);
  \draw (2.32,0) -- (2.58,0);
  \draw (2.32,4) -- (2.58,4);
  \node at (2.85,2) {$4$};
  \node at (1,-1.35) {正视图};

  % ---------- 侧视图 ----------
  \begin{scope}[xshift=4.1cm]
    \draw (0,0) rectangle (2,4);
    \draw (1,3) circle (1);
    \draw[<->] (0,-0.45) -- (2,-0.45);
    \draw (0,-0.58) -- (0,-0.32);
    \draw (2,-0.58) -- (2,-0.32);
    \node at (1,-0.85) {$2$};
    \node at (1,-1.35) {侧视图};
  \end{scope}

  % ---------- 俯视图 ----------
  \begin{scope}[yshift=-3.55cm]
    \draw (1,0) circle (1);
    \node at (1,-1.30) {俯视图};
  \end{scope}
\end{tikzpicture}
